\documentclass[11pt]{article}
\usepackage{amsmath,amssymb,amsthm,mathrsfs}
\usepackage[margin=1in]{geometry}
\usepackage[hidelinks]{hyperref}
\newcommand{\DD}{d'd''}
\newcommand{\MA}{\operatorname{MA}}
\newcommand{\C}{\mathcal C}
\newcommand{\G}{\mathcal G}
\newcommand{\GN}{\mathcal G_{\mathrm{net}}}
\newcommand{\GS}{\mathcal G_{\mathrm{seq}}}
\newcommand{\RR}{\mathbb R}
\newcommand{\HH}{\mathscr H}
\newtheorem{theorem}{Theorem}
\newtheorem{lemma}[theorem]{Lemma}
\newtheorem{corollary}[theorem]{Corollary}
\newtheorem{proposition}[theorem]{Proposition}
\theoremstyle{remark}\newtheorem{remark}[theorem]{Remark}
\title{Energy-bounded decreasing completions on a nonmetrizable analytic space\\Actual nets, countable flattening, and finite Shilov coordinates}
\author{Research working paper 406; unreviewed}
\date{9 October 2026}
\begin{document}
\maketitle
\begin{abstract}
We give an actual continuous presentation for the infimum of an
energy-bounded decreasing net of continuously generated weights on
a compact Hausdorff space. Its energy is canonical and equals the
decreasing limit of the outer energies. In the sequential case a
countable diagonal of the supplied presentations suffices, without
a countable dense set or a metrizability assumption. Finite full
bounded-energy dual measures give genuine integrability and
$L^1$ convergence for every presentation net, using Radon
regularity rather than a false arbitrary-net monotone convergence
principle for measurable functions. On the actual analytic surface,
bounded atomic tests and classical Shilov reconstruction extend to
the net-generated class. These results do not identify the sequential
and net classes or prove representability of arbitrary product
limits. Original P16.7 remains open. All new claims are unreviewed.
\end{abstract}

\section{The class being completed}
For the first six sections let $X$ be a nonempty compact Hausdorff
space. Let $\C\subset C(X,\RR)$ contain every constant and be stable
under addition of constants. Suppose $E:\C\to\RR$ is monotone for
the pointwise order and satisfies
\[
 E(0)=0,\qquad E(a+c)=E(a)+Mc,\qquad M>0.
\]
No lattice closure, envelopes, measure operator, metrizability,
countable cofinality, or continuity of $E$ is assumed in these
abstract arguments. These hypotheses hold for the actual continuous
CLD class and raw energy on the analytic surface described in the
appendix, where $M=2$.

A directed index set is nonempty and every finite set of indices
has an upper bound. Write $a_i\downarrow u$ when $a_i\in\C$ is
decreasing in this order and $u(x)=\inf_i a_i(x)$ for every $x$.
Define $\GN$ by such actual presentations with
$\inf_iE(a_i)>-\infty$. Define $\GS$ by presentations indexed by
$\mathbb N$ with the same bound. Thus $\GS\subset\GN$; equality
is not part of the definition or a conclusion below. The weights
are actual functions $X\to\RR\cup\{-\infty\}$, not measure
equivalence classes. Each is upper semicontinuous and is bounded
above after restricting to the tail of any fixed presentation index.
Restricting to a tail changes neither the pointwise infimum nor
the energy infimum, by directedness.

\section{Compact strict domination and canonical energy}
\begin{lemma}\label{strict}
Let $f_i$ be a decreasing net of upper semicontinuous functions,
let $f=\inf_i f_i$, and let $a$ be continuous with $f<a$ everywhere.
There is an index $i_0$ such that $f_i<a$ everywhere for every
$i\ge i_0$. If $f_{i_0}<a$, there is a real $\delta>0$ such that
$f_{i_0}<a-\delta$ everywhere.
\end{lemma}
\begin{proof}
The ordinary open sets $\{f_i<a\}$ cover $X$. Choose a finite
subcover and an upper bound of its indices. Decrease then gives
the first assertion. For the second assertion the open sets
$\{f_{i_0}<a-1/n\}$, $n\ge1$, cover $X$ and are increasing;
compactness gives one of these sets equal to $X$.
This argument includes poles and uses no countable cofinality
of the original directed index set.
\end{proof}

\begin{theorem}\label{canonical}
For $u\in\GN$, every continuous decreasing presentation has the
same finite energy limit, denoted $E(u)$, and
\begin{equation}\label{majorant}
 E(u)=\inf\{E(a):a\in\C,\ a>u\text{ everywhere}\}.
\end{equation}
If $u,v\in\GN$ and $u\le v$, then $E(u)\le E(v)$.
For every finite constant $c$, $u+c\in\GN$ and
$E(u+c)=E(u)+Mc$.
No element of $\GN$ is identically $-\infty$; its supremum is a
finite real number, and $E(u)\le M\sup_Xu$.
\end{theorem}
\begin{proof}
Fix a presentation $a_i\downarrow u$ and set
$e=\inf_iE(a_i)\in\RR$. If $a\in\C$ is strictly above $u$,
Lemma\ref{strict} gives some $a_i<a$ globally, so $e\le E(a)$.
Conversely $a_i+\epsilon>u$ is a permitted majorant, and
$E(a_i+\epsilon)=E(a_i)+M\epsilon$. Infimizing over $i$ and
$\epsilon>0$ proves\eqref{majorant}. In particular the right side
is independent of the presentation. If another presentation has
no stated lower bound, the same strict-domination comparison
forces its energy infimum to equal $e$, so it has the bound as well.
Monotonicity follows because every strict continuous majorant of
$v$ is also a majorant of $u$.
Shifting an actual presentation by $c$ proves the stated constant
rule and membership.

If $u\equiv-\infty$, compact strict domination eventually places
the presentation below each constant $-N$, forcing $e\le-MN$ for
every $N$. This contradicts finiteness. The first continuous term
of a tail gives a finite upper bound, while a finite value of $u$
gives a finite lower bound for its supremum. Constants
$\sup_Xu+\epsilon$ in\eqref{majorant} give the last inequality.
\end{proof}

The assertion about an arbitrary second presentation can also be
seen directly: for fixed $a_i+\epsilon$, a term of the second
presentation lies below it, giving its energy infimum at most
$e$; for every term $b_j$ of the second presentation,
$b_j+\epsilon>u$, and strict domination by the first presentation
gives $e\le E(b_j)+M\epsilon$. Let $\epsilon$ decrease to zero
and infimize over $j$.

\section{Actual flattening of an energy-bounded decreasing net}
\begin{theorem}\label{netflat}
Let $u_\gamma\in\GN$ be a decreasing net with
$E(u_\gamma)\ge L$ for a real $L$. Then its pointwise infimum
$u=\inf_\gamma u_\gamma$ belongs to $\GN$ and
\[
 E(u)=\inf_\gamma E(u_\gamma)\ge L.
\]
An actual continuous presentation is the net of all members of
\begin{equation}\label{D}
 \mathscr D_L(u)=\{a\in\C:a>u\text{ everywhere},\ E(a)\ge L\},
\end{equation}
ordered by reverse pointwise order. This set is nonempty,
downward filtered, and has pointwise infimum $u$.
\end{theorem}
\begin{proof}
Fix any outer index $\gamma$ and a continuous presentation
$c_i\downarrow u_\gamma$. Every $c_i$ has energy at least
$E(u_\gamma)\ge L$. For $\epsilon>0$, $c_i+\epsilon>u_\gamma\ge u$
and has energy at least $L$, so\eqref{D} is nonempty.

Given $a,b\in\mathscr D_L(u)$, their ordinary continuous minimum
$m=\min(a,b)$ is strictly above $u$. We do not assume $m\in\C$.
Apply Lemma\ref{strict} to the outer USC net to choose an index
$\gamma$ and a number $\delta>0$ with $u_\gamma<m-\delta$ globally.
Apply the same lemma to a continuous presentation $c_i$ of this
$u_\gamma$ to choose $c_i<m-\delta/2$ everywhere. Pick
$0<\epsilon<\delta/2$. Then
\[
 c_i+\epsilon\in\mathscr D_L(u),\qquad
 c_i+\epsilon<a,\qquad c_i+\epsilon<b.
\]
Thus the set is downward filtered even though $\C$ need not be
closed under minima.

Every member of\eqref{D} is above $u$. Fix $x\in X$ and a real
$T>u(x)$, where any real $T$ is permitted at a pole. Choose
$\gamma$ with $u_\gamma(x)<T$, then a term $c_i$ of its
presentation with $c_i(x)<T$. A sufficiently small positive
$\epsilon$ leaves $(c_i+\epsilon)(x)<T$ while making it a member
of\eqref{D}. Hence the infimum is exactly $u$ at every point.
The net so obtained has energy bounded below by $L$, proving
actual membership in $\GN$ before its energy is used.

Monotonicity in Theorem\ref{canonical} gives
$E(u)\le E(u_\gamma)$ for every $\gamma$. Conversely, for any
$a\in\C$ strictly above $u$, compact strict domination gives
$u_\gamma<a$ everywhere for some $\gamma$.
By\eqref{majorant}, $E(u_\gamma)\le E(a)$. Taking the infimum
over $\gamma$ and then over all such $a$ gives
$\inf_\gamma E(u_\gamma)\le E(u)$. This proves equality.
\end{proof}

\begin{remark}
The index set\eqref{D} is an actual set of finite continuous
functions. It is not a presumed cofinal family of algebraic models,
an abstract functional completion, or a prescribed sequence.
In particular the theorem supplies decreasing-net closure of the
net-generated class at a common finite energy bound without
enlarging its definition by an iterated completion.
\end{remark}

\section{A countable diagonal using the supplied presentations}
\begin{theorem}\label{seqflat}
Suppose $u_n\in\GS$ decreases and $E(u_n)\ge L$.
Then $u=\inf_nu_n\in\GS$ and $E(u)=\inf_nE(u_n)$.
This conclusion uses no countable dense set of points of $X$.
\end{theorem}
\begin{proof}
Choose each supplied presentation $a_{n,k}\downarrow u_n$.
All its terms satisfy $E(a_{n,k})\ge L$.
Start with $v_1=a_{1,1}+1$, which is strictly above $u_1$.
Given $v_{n-1}>u_{n-1}$, the finite collection
\[
 v_{n-1},\qquad a_{i,k}+2^{-n}\quad(1\le i,k\le n)
\]
consists of continuous functions strictly above $u_n$.
Let $m_n$ be their ordinary continuous minimum. Compactness gives
a number $\delta_n>0$ with $u_n<m_n-\delta_n$ everywhere.
Compact strict domination chooses $k(n)$ with
$a_{n,k(n)}<m_n-\delta_n/2$ everywhere. Pick
$0<\epsilon_n<\min(\delta_n/2,2^{-n})$ and set
$v_n=a_{n,k(n)}+\epsilon_n$. Then $v_n\in\C$,
$v_n>u_n$, $v_n<v_{n-1}$ and $E(v_n)\ge L$.

Write $v=\inf_n v_n$. It is at least $u$.
For fixed $i,k$, every sufficiently large $n$ gives
$v_n<a_{i,k}+2^{-n}$, hence $v\le a_{i,k}$.
Infimizing first over $k$ and then over $i$ gives $v\le u$.
Thus this actual sequence presents $u$ with the required energy
bound. The energy equality follows from Theorem\ref{netflat},
which applies to the same outer sequence, and from canonical energy.
\end{proof}

\section{Radon integrals for monotone continuous nets}
Arbitrary directed monotone convergence for measurable functions
is false: indicators of finite subsets of $[0,1]$, directed by
inclusion, have Lebesgue integrals zero but supremum one.
The following topological version is the one needed here.

\begin{lemma}\label{radon}
Let $\mu$ be a finite positive Radon measure on compact Hausdorff
$X$, and let $g_i$ be an increasing net of nonnegative continuous
functions. For $g=\sup_i g_i$, possibly infinite,
\[
 \int g\,d\mu=\sup_i\int g_i\,d\mu.
\]
\end{lemma}
\begin{proof}
The function $g$ is lower semicontinuous and therefore Borel.
Radon regularity gives
\begin{equation}\label{lscint}
 \int g\,d\mu=
 \sup\{\int h\,d\mu:h\in C(X,\RR),\ 0\le h\le g\}.
\end{equation}
Here is a direct verification of the nontrivial inequality.
For a finite step $\delta>0$ and cutoff $N$, choose compact sets
$K_k\subset\{g>k\delta\}$, $1\le k\le N$, whose measures
approximate these open-set measures from below. Normality of
compact Hausdorff $X$ gives $h_k\in C(X,[0,1])$, equal to one
on $K_k$ and with support inside $\{g>k\delta\}$.
Then $h=\delta\sum_{k=1}^N h_k$ is continuous and at most $g$:
only indices $k$ with $k\delta<g(x)$ can contribute at $x$.
Its integral approximates the finite lower layer-cake sums.
Letting the compact-set errors tend to zero, then refining
$\delta$ and increasing the cutoff, proves\eqref{lscint},
including the infinite-integral case.

For $0\le h\le g$ and $\epsilon>0$, the open sets
$\{g_i>h-\epsilon\}$ cover $X$. A finite subcover and
directedness give a single $i$ with $g_i>h-\epsilon$ globally.
Consequently $\sup_i\int g_i\,d\mu\ge\int h\,d\mu-
\epsilon\mu(X)$. Use\eqref{lscint} and let $\epsilon$ tend to
zero. The opposite inequality is immediate from $g_i\le g$.
\end{proof}

\section{The full bounded dual controls actual net presentations}
For a finite positive Radon measure define
\[
 E^*(\mu)=\sup_{a\in\C}\left(E(a)-\int a\,d\mu\right).
\]
Constants show that $E^*(\mu)<\infty$ forces $\mu(X)=M$;
the zero competitor gives $E^*(\mu)\ge0$.

\begin{theorem}\label{dualmoment}
If $E^*(\mu)<\infty$ and $u\in\GN$ with $u\le0$, then
\[
 \int(-u)\,d\mu\le E^*(\mu)-E(u)<\infty.
\]
Every decreasing continuous presentation converges to $u$ in
$L^1(\mu)$ as an actual net. For a general $u$, the same assertion
holds after a finite constant shift, so poles are $\mu$-null.
\end{theorem}
\begin{proof}
Fix a presentation and a tail bounded above by a constant $B$.
The continuous nonnegative functions $B-a_i$ increase to $B-u$.
For each $i$, the definition of the dual gives
\[
 \int(B-a_i)\,d\mu\le MB+E^*(\mu)-E(a_i).
\]
Lemma\ref{radon} and the canonical energy limit give
$\int(B-u)\,d\mu\le MB+E^*(\mu)-E(u)$.
This establishes integrability before using an integral of $u$.
Subtracting $MB$ proves the displayed bound. Since $a_i\ge u$,
\[
 \|a_i-u\|_{L^1(\mu)}=\int a_i\,d\mu-\int u\,d\mu
 \longrightarrow0
\]
by the same integral limit. Finite supremum and the constant rule
in Theorem\ref{canonical} justify the stated shift.
\end{proof}

\section{Application to actual analytic weights}
Now specialize to the actual surface
\[
 K=\mathbb C((s)),\quad |s|=e^{-1/8},\quad t=s^8,
 \qquad X=(\mathbb G_m^2)^{\mathrm{an}}_K/
          \langle(t,2),(3,t)\rangle.
\]
Use the all-stage continuous CLD sheaf $\mathscr P$, local
$\rho=|\tau|^2/2$, and $\theta=\DD\rho$ in the appendix.
Its bounded Green calculus proves the energy hypotheses and
\[
 E^*(\MA(w))=E(w)-\int w\MA(w)<\infty\qquad(w\in\C).
\]
The appendix reproduces the complete atomic construction: at
each point $x$ of graded dimension two there is an actual bounded
$w_x\in\C$ with
\[
 \MA(w_x)\ge\alpha_x\delta_x,\qquad\alpha_x>0.
\]
Its raw calibrated coefficient and finite flat map are established
there without a change of ground field. This is an unreviewed
working foundation whose full proof is supplied, not a new
independent correctness grant or a consequence of compilation.

\begin{corollary}\label{netpoints}
For every normalized $u\in\GN$ with $\sup_Xu=0$ and
$E(u)\ge-C$, all graded-dimension-two values are finite and
\[
 0\ge u(x)\ge-
 \frac{E^*(\MA(w_x))+C}{\alpha_x}.
\]
\end{corollary}
\begin{proof}
Apply Theorem\ref{dualmoment} to the whole mass-$M$ measure
$\MA(w_x)$. Positivity gives
$\alpha_x(-u(x))\le\int(-u)\MA(w_x)\le E^*(\MA(w_x))+C$.
No finite-dual assertion for the isolated component $\delta_x$
is used.
\end{proof}

The all-stage CLD-to-classical recognition in the appendix applies
to continuous presentation terms. The literal GRv2
Proposition4.2(4) explicitly permits a decreasing \emph{net} of
classically psh functions. Thus every local weight $\rho+u$ for
$u\in\GN$ is actual USC classical psh, including at poles.
Theorem6.4 gives its maximum on each affinoid Shilov boundary.
For an arbitrarily small affinoid neighborhood $A$ of $y$ in
one quotient chart, it follows that
\[
 (\rho+u)(y)=\inf_{A\ni y\,\mathrm{neighborhood}}
 \max_{\xi\in\Gamma_A}(\rho+u)(\xi).
\]
One inequality is the maximum principle; the other follows from
USC by shrinking $A$, using every finite upper threshold when
$u(y)=-\infty$. The union $\mathscr S$ of these Shilov boundaries
therefore determines actual net-generated weights at every point.
Every such Shilov point has graded dimension two, by the exact
general-radius GRv2 Lemma6.3 and the Abhyankar inequality retained
in the appendix. Corollary\ref{netpoints} bounds all its values.

\section{What this closure theorem supplies for compactness}
The evaluation map on $\mathscr S$ embeds the normalized
energy sublevel of $\GN$ as a set into
\[
 \prod_{\xi\in\mathscr S}[-L_\xi(C),0],\qquad
 L_\xi(C)=\frac{E^*(\MA(w_\xi))+C}{\alpha_\xi}.
\]
The ambient product is compact. Theorems\ref{netflat} and
\ref{seqflat} resolve actual representability of a
\emph{decreasing} outer family with a common finite energy bound.
They do not resolve representability of an arbitrary convergent
subnet of this product. A product limit need not be a decreasing
limit of supplied weights, and a continuous positive majorant
directed family cannot be presumed to exist for a merely specified
coordinate tuple. No closed-image or energy-semicontinuity theorem
for arbitrary coordinate limits is claimed.

In particular we have not proved $\GS=\GN$, a countable cofinal
presentation for arbitrary nets, upper-envelope positivity or
regularity, Monge--Ampere existence for arbitrary finite-dual
measures, or full original P16.7. A Radon measure with prescribed
properties cannot be replaced by countably many point tests without
proof. These remain separate mathematical tasks.

\section{Reading scope and reproducibility}
The new abstract proofs above are supplied in full. They use
ordinary compact Hausdorff normality, finite-cover directedness,
Radon inner regularity on open sets, Urysohn functions, and the
ordinary scalar layer-cake integral formula at the stated finite
measure hypotheses. They use no arbitrary-measurable-net monotone
convergence theorem.

The complete immediate GRv2 Proposition4.2 statement and proof
were physically read with their surrounding LF1116--1198 context.
Its decreasing-net assertion reduces to Proposition3.5(5), whose
recursive whole proof is not freshly claimed read here.
The smooth recognition statement is the unnumbered argument of
Section1.2, not a theorem numbered1.2. The whole exact affinoid
maximum statement and immediate proof of Theorem6.4, and the
whole Lemma6.3 Shilov-dimension argument, retain the appendix's
qualified primary reading scope. Deeper Temkin, Ducros, Berkovich,
intersection and calibration results remain named published
foundations, not an invented new recursive audit.

The appendix is the complete literal mathematical body of the
current citation-corrected paper402. Exact identity is retention,
not a second independent proof reading. Its historical residual
language describes that paper's scope at writing and grants no
new status. All primary bodies, literal source units, pixels and
book contexts remain LOCALONLY. All new claims are UNREVIEWED;
original closure and counts remain unchanged at16/42 solved,
26OPEN and97accepted scoped papers. Native/generated programs
are DATA and are not executed or imported. The protected P16.4
editor remains OPEN UNCHANGED; this new paper is not hot feedback
to the active frozen campaigns.

\appendix
\section{Complete point-test and bounded-energy foundation from paper402}
The following full mathematical body is retained literally.
Its section structure and labels are preserved; its actual
notation is shared with the main text.
\section{The actual analytic object and proof foundations}
Put $K=\mathbb C((s))$, $|s|=e^{-1/8}$, $t=s^8$,
$Y=(\mathbb G_m^2)^{\mathrm{an}}_K$, and
\[
 X=Y/\langle(t,2),(3,t)\rangle,\qquad
 \tau=(-\log|z_1|,-\log|z_2|).
\]
The translations on $\tau$ form $\mathbb Z^2$. Ordinary open tropical
boxes of width less than one give injective quotient charts. We use
the previously reconstructed analytic quotient: compact, separated,
good, boundaryless, connected, smooth, and pure of dimension two.
The quotient construction and its nonalgebraization proof are not
reproved here and are not consequences of the present point test.
On a chart let $\rho=|\tau|^2/2$ and $\theta=\DD\rho$.
Changes of reference are affine unit logarithms. With the identity
torus's saturated lattice and degree one,
\[
 M=\int_X\theta^2=2.
\]
This is the raw CLD convention; no probability normalization is made.

Let $\mathscr P$ be the minimal continuous CLD psh sheaf containing
the ordinary smooth psh functions and closed under local uniform
limits. All stages of this sheaf are retained. Set
\[
 \C=\{a\in C(X,\RR):\rho+a\in\mathscr P\text{ locally}\},\quad
 \theta_a=\theta+\DD a,\quad B(a,b)=\theta_a\wedge\theta_b.
\]
We use the actual recursive continuous products. Their positivity,
symmetry, closedness, locality, local uniform continuity, and
invariance under smooth $\DD$-null reference changes are precisely
CLDv2 8.5.1, 8.5.3, 8.5.5--6, 8.7.1--4 and 8.8.4--5, at the stated
continuous psh/dpsh and positive closed-current hypotheses. Positive
top currents are Radon measures by 7.2.8 on these locally
differentially separated charts. Smooth cutoffs, partitions, scalar
approximation and boundaryless Stokes are the conventional analytic
foundations. Arbitrary continuous-multiplier closedness in 8.6.7 is
not a premise. No first-stage identity for $\mathscr P$ is assumed.

Constants, convex combinations and finite maxima preserve $\C$.
For a smooth slot, convex regularized maxima are smooth psh.
Uniform continuity and sheaf minimality extend the operation one
slot at a time to all stages. This also proves ordinary convex
composition with finite continuous psh slots by the same argument.

\section{Bounded Green calculus and the full bounded dual}
Here we give the extra energy argument rather than infer it from
a positivity label. For a positive closed $(1,1)$ current $T$,
complete compact smooth $(1,0)$ forms modulo the null space of
\[
 (a,b)\longmapsto\int a\wedge Jb\wedge T.
\]
Use a nonnegative compact smooth cutoff $\chi$ for the corresponding
local seminorm $Q_{T,\chi}$. Let $\HH_T$ consist of continuous
functions locally uniformly approximable by smooth functions whose
gradients converge strongly in these seminorms. The gradient is
unique: for a uniformly zero sequence $f_j$ with a strong gradient
limit and a compact smooth form $b$, closedness gives
\[
 \int d'f_j\wedge Jb\wedge T=-\int f_jd'(Jb)\wedge T\longrightarrow0.
\]
The right-hand coefficient measure is finite on the fixed compact
support; smooth Hilbert tests are dense. The limit gradient is zero.
Partitions patch local approximants by the Leibniz rule.
Bounded continuous coefficients act by scalar approximation and the
bound $\|fa\|\le\|f\|_\infty\|a\|$. Products and cutoffs have their
usual gradient rules.

If $f\in\HH_T$ and $L_Tf=\DD(fT)$ is signed Radon, smooth
approximation first in $f$, then uniformly and strongly in a
compactly supported $g$, yields
\[
 \int gL_Tf=-Q_T(f,g),\qquad
 Q_{T,\chi}(w,w)=\frac12\int w^2\DD\chi\wedge T-\int\chi wL_Tw.
\]
The sheaf of $f\in\HH_T$ with $L_Tf\ge0$ contains smooth psh
functions and is uniformly closed. Indeed take nested cutoffs
$\chi,\eta$, $\eta=1$ near $\operatorname{supp}\chi$,
a common bound $|f_j|\le B$ on the whole support of $\eta$, and
$\epsilon=\|f_j-f_k\|$ there. Positivity and the preceding identity give
\[
 0\le Q_{T,\chi}(f_j-f_k,f_j-f_k)
 \le\frac{\epsilon^2}{2}\|\DD\chi\wedge T\|_{\mathrm{TV}}
 +2B\epsilon\|\chi\|_\infty\|\DD\eta\wedge T\|_{\mathrm{TV}}.
\]
Choose smooth approximants to $f_j$ with uniform and smaller-patch
gradient errors below $1/j$. The estimate proves strong convergence;
distributional positivity survives. Minimality places every stage
of $\mathscr P$ in this sheaf. Its vector span includes continuous
dpsh functions. Finite patching on compact $X$ proves
\[
 \int f\DD(gT)=\int g\DD(fT)=-Q_T(f,g),\qquad Q_T(f,f)\ge0.
\]
In particular all the following computations hold at every
continuous stage. The exact terms in the bounded product expansion
have zero total integral, so every $B(a,b)$ has mass $M=2$.

Define
\[
 E(a)=\frac13\int_X a[\theta^2+\theta_a\wedge\theta+\theta_a^2].
\]
For $a_t=(1-t)a+tb$, the cubic expansion and Green give
\[
 \frac{d}{dt}E(a_t)=\int(b-a)\MA(a_t),\qquad
 \frac{d^2}{dt^2}E(a_t)=-2Q_{\theta_{a_t}}(b-a,b-a)\le0.
\]
Thus $E$ is concave and monotone, $E(0)=0$, and
$E(a+c)=E(a)+Mc$. Its tangent inequality is
\[
 E(a)\le E(w)+\int(a-w)\MA(w)\qquad(a,w\in\C).
\]
For a positive Radon measure $\mu$ define the \emph{full} bounded dual
\[
 E^*(\mu)=\sup_{a\in\C}\left(E(a)-\int a\,d\mu\right).
\]
Constants force $\mu(X)=M$ whenever this quantity is finite.
The tangent inequality, with equality at $a=w$, proves the exact
formula
\begin{equation}\label{dual}
 E^*(\MA(w))=E(w)-\int w\MA(w)<\infty
 \quad\text{for every }w\in\C.
\end{equation}
It is nonnegative, since $a=0$ is one of the competitors.

\section{The actual generated class and a moment bound}
Let $\G$ consist of actual pointwise limits
$a_j\downarrow u$ on $X$, where $a_j\in\C$ and
$\inf_jE(a_j)>-\infty$. These limits are USC.
They are not abstract functionals or equivalence classes modulo
null sets. Define $E(u)=\lim_jE(a_j)$.
This is independent of every continuous decreasing presentation.
To see this, if $b_k\downarrow u$ and $a$ is continuous with $u<a$
everywhere, the increasing ordinary open sets $\{b_k<a\}$ cover
compact $X$. One index therefore works everywhere. Apply this to
$a=a_j+\epsilon$, use monotonicity and the constant rule, and
then reverse the roles of the presentations. The limits agree.
In particular $E(u)>-\infty$ for every presentation of a member
of $\G$. The identically $-\infty$ limit is excluded: compact
strict domination would eventually put a presentation below every
negative constant, contradicting its energy lower bound.

\begin{lemma}\label{moment}
If $u\in\G$, $u\le0$, and $E^*(\mu)<\infty$ for a positive
Radon measure $\mu$, then
\[
 \int(-u)\,d\mu\le E^*(\mu)-E(u)<\infty.
\]
\end{lemma}
\begin{proof}
Every term of a presentation satisfies
$\int(-a_j)\,d\mu\le E^*(\mu)-E(a_j)$.
The continuous first term bounds all terms above. Apply monotone
convergence to that common bound minus $a_j$ and pass to the limit.
This proves the claim even at poles and does not presume
integrability before the bound. The same argument proves
$L^1(\mu)$ convergence of the whole presentation.
\end{proof}

\section{Laurent Abhyankar moments without changing the field}
Write $d_K(x)=\operatorname{trdeg}_{\widetilde K}
\widetilde{\mathscr H(x)}$, with graded residue fields.
Let $\mathscr A$ be the set of points of $X$ with $d_K(x)=2$.
A lift of such a point in $Y$ will also be denoted $x$.

\begin{lemma}\label{moments}
There are two nonzero Laurent polynomials $P_1,P_2$ over $K$
whose initial forms at $x$ are algebraically independent over
$\widetilde K$. Put $r_i=|P_i(x)|>0$. Then
\[
 (P_1,P_2)(x)=\eta_{(r_1,r_2)}
\]
in $(\mathbb A_K^2)^{\mathrm{an}}$. The map is finite and flat at
$x$, and admits a local finite flat representative with a positive
constant degree $D$ and with $x$ the only point above this Gauss point.
\end{lemma}
\begin{proof}
The support of $x$ on the Laurent polynomial ring is zero.
Otherwise it lies on a proper algebraic subvariety of dimension
at most one, contradicting $d_K(x)=2$ and the Abhyankar inequality.
Thus $\mathscr H(x)$ is the completion of the valued field
$K(z_1,z_2)$. Completion does not add graded residue elements:
approximate a nonzero element with error strictly less than its
norm. Its initial form agrees with that of the approximant.
Consequently the initial forms of all nonzero Laurent polynomials
generate the graded residue \emph{field}, allowing ratios.
A maximal algebraically independent subfamily of these homogeneous
generators has cardinal two. Choosing its members gives $P_1,P_2$.

For every polynomial $\sum_I c_IP^I$, the initial terms of maximal
norm cannot cancel by graded algebraic independence. Its norm is
$\max_I|c_I|r^I$. This is exactly the Gauss seminorm, for arbitrary
real $r_i>0$; no scaling of $r_i$ into $|K^\times|$ is required.
This is also the Abhyankar-basis criterion of CLD 13.8.5.

The torus is boundaryless, hence $x$ is in the relative interior
of the map to $\mathbb A^2$. The target is reduced and the Gauss
point has graded dimension two. CLD 11.6.5(c) makes the map finite
and flat at $x$. Its 11.6.3 local representative can be chosen
with a single point of that fibre. The rank is a positive integer
$D$. After shrinking the target inside its ordinary interior,
the source is an ordinary open neighborhood of $x$, and the finite
flat degree remains $D$ there.
\end{proof}

\section{The actual atomic mass and its global test}
We spell out the local calibrated number.
On a real two-plane set $g(t)=\max(0,t_1,t_2)$ and
$g_\epsilon=\epsilon\log(1+e^{t_1/\epsilon}+e^{t_2/\epsilon})$.
For $g_1$, write $p_i=e^{t_i}/(1+e^{t_1}+e^{t_2})$ and
$p_0=1-p_1-p_2$. The Hessian determinant is $p_0p_1p_2$.
The gradient maps the plane diffeomorphically to
$\{p_1,p_2>0:p_1+p_2<1\}$ with that Jacobian.
The raw superform rule is $2!\det(\operatorname{Hess}g_1)\,dt_1dt_2$.
Its integral is $2!\,(1/2)=1$.
Dilation and dominated convergence concentrate this finite measure
at the origin as $\epsilon\downarrow0$, so
\[
 (\DD g)^2=\delta_0.
\]
This is the explicit case of CLD 1.7.1 and 1.7.6--8, including
their factor $2!$.

Near $x$, both $P_i$ are invertible. Apply the pleasant moment
$(P_1,P_2)$ there. Finite flatness and the single Gauss fibre imply
full tropical dimension and that the calibrated tropical cone at
$x$ is the whole real plane with multiplicity $D$.
Indeed the map is open, so the Gauss section of a neighborhood of
the target gives a neighborhood in the tropical plane. The explicit
calibration formula CLD 14.5.2--3 makes its image calibration
$D$ times the standard one. Locality removes other fibre points,
since there are none in this representative. More explicitly,
smooth log-sum-exp approximants have top measures strongly
supported on the inverse image of the target Gauss section.
The calibrated integration formula pushes them to $D$ times the
corresponding real-plane measures near the target tropical value.
Those measures concentrate there with mass one. Finite properness
and the single Gauss fibre imply that every inverse Gauss point
over a sufficiently small tropical neighborhood lies in any
prescribed ordinary neighborhood of $x$: otherwise compactness
would give another point of that fibre. Continuous tests therefore
converge to $D$ times evaluation at $x$. Local uniform continuity
of the actual product identifies the limit with the log-maximum
product. This reconstructs the atomic calculation 8.9.2--4 and yields
\begin{equation}\label{localatom}
 \left(\DD\max(0,\log(|P_1|/r_1),\log(|P_2|/r_2))\right)^2
 =D\delta_x
 \quad\text{on some ordinary neighborhood of }x.
\end{equation}
Formula \eqref{localatom} uses the normalized
canonical calibration over the unchanged field $K$ and the actual
finite flat degree. No arbitrary-current pullback is invoked.
CLD 18.6.2 gives the equivalent Gauss-fibre description, but its
relative-intersection proof is not needed for this argument.

\begin{theorem}\label{atom}
For every $x\in\mathscr A$ there is $w_x\in\C$ and
$\alpha_x>0$ such that the actual positive Radon measure
$\MA(w_x)$ satisfies
\[
 \MA(w_x)(\{x\})=\alpha_x.
\]
One can take $\alpha_x=\varepsilon^2D$ in the construction below.
\end{theorem}
\begin{proof}
Choose an injective lift chart around the tropical class
$a=\tau(x)$ and a small Euclidean tropical ball of radius $R$ there.
Its inverse image has compact closure inside a polyannulus.
Set $\rho_a=|\tau-a|^2/2$; its difference from $\rho$ is
smooth affine and $\DD$-null.
The function
\[
 \psi=\max(0,\log(|P_1|/r_1),\log(|P_2|/r_2))
\]
is finite, continuous, nonnegative and psh on this whole chart,
including zeros of the polynomials. Here is a direct check at all
continuous stages. Choose a smooth convex $\chi$ equal to zero for
$v\le0$, with $0\le\chi'\le1$, and equal to $v-c$ for $v\ge1$,
where $0<c<1$. For $N>3$ the finite smooth psh function
\[
 p_{i,N}=\chi(\log(|P_i|/r_i)+N)-N+c
\]
is locally constant at zeros. Where it can be positive it equals
$\log(|P_i|/r_i)$, so $\psi=\max(0,p_{1,N},p_{2,N})$ exactly.
Smooth convex regularized maxima place it in $\mathscr P$.

Take $0<\delta<R^2/4$ and choose $\varepsilon>0$ so small that
$\delta+\varepsilon\sup\psi<R^2/2$ on the compact closure.
Then
\[
 w_x=\max(\rho_a,\delta+\varepsilon\psi)-\rho_a
\]
vanishes on a tropical boundary collar. Extend it by zero to $X$.
It belongs to $\C$: regularize the inner finite maximum, and then
use a smooth convex regularized outer maximum which agrees exactly
with the larger slot when the gap exceeds its regularization
width. The fixed collar gap ensures that each sufficiently small
regularization extends by zero smoothly and $\theta$-psh.
Their uniform limit is the displayed function.

On the fixed inner neighborhood $\rho_a<\delta/2$ we have
$\theta_{w_x}=\varepsilon\DD\psi$. Shrink further to the
neighborhood in \eqref{localatom}. Locality of the actual recursive
product gives
$\MA(w_x)=\varepsilon^2D\delta_x$ there.
Outside it the global product remains positive and has raw mass
$M=2$. Thus its atom at $x$ is precisely the positive number
$\alpha_x=\varepsilon^2D$.
\end{proof}

\section{Finiteness, sublevel bounds, and the precise compactness limit}
\begin{corollary}\label{finite}
For every $u\in\G$ and every $x\in\mathscr A$, $u(x)$ is finite.
More precisely, if $\sup_Xu=0$ and $E(u)\ge-C$, then
\[
 -\frac{E^*(\MA(w_x))+C}{\alpha_x}\le u(x)\le0 .
\]
\end{corollary}
\begin{proof}
Use \eqref{dual} and Lemma \ref{moment} for the \emph{whole}
positive measure $\mu_x=\MA(w_x)$.
Since $-u\ge0$ and $\mu_x$ contains the positive atom
$\alpha_x\delta_x$,
\[
 \alpha_x(-u(x))\le\int(-u)\,d\mu_x
 \le E^*(\mu_x)-E(u)\le E^*(\mu_x)+C.
\]
This excludes a pole at $x$ and proves the bound.
An arbitrary $u$ can be shifted by its finite supremum.
A member of $\G$ has a finite upper bound and is not identically
$-\infty$, so that shift is legitimate.
\end{proof}

Every Shilov point of a pure-dimensional two-dimensional affinoid
domain in a local chart has $d_K=2$. At the exact general-radius
hypotheses, GRv2 Lemma6.3(3)--(4) expresses these Shilov points as
the strata with two graded algebraically independent generators;
the Abhyankar inequality gives the reverse bound.
The entire lemma and its immediate proof were read, including its
Temkin graded-reduction and Ducros fibre-dimension foundations.
These deeper theorems are named published foundations, not freshly
recursively reproved here. Hence Corollary \ref{finite} applies to
every such Shilov point, not only to the reference torus skeleton.
On a fixed finite Shilov boundary the lower bound can be chosen
uniformly by taking the maximum of the finitely many displayed
constants.

For clarity, actual local weights are classically psh at all
stages. The published smooth CLD-to-classical bridge is the
unnumbered statement and argument in GRv2 Section1.2.
Its Proposition4.2 gives sheafiness and decreasing-net
closure. A uniform limit of finite continuous classical psh
functions is classical psh: select errors at most $2^{-j}$ and add
$3\cdot2^{-j}$, obtaining a decreasing sequence with the same limit.
Sheaf minimality then recognizes all of $\mathscr P$, and a
continuous decreasing presentation recognizes $\rho+u$ at poles.
GRv2 Theorem6.4 therefore gives the affinoid Shilov maximum
principle for these actual USC weights.
For all arbitrarily small affinoid neighborhoods $A$ in one chart,
\[
 (\rho+u)(y)=\inf_{A\ni y\ \mathrm{neighborhood}}
                \max_{\xi\in\Gamma_A}(\rho+u)(\xi).
\]
The maximum principle gives one inequality. USC gives the other
by every finite upper bound at a pole and by
$(\rho+u)(y)+\eta$ at a finite value. Thus values on the union
$\mathscr S$ of \emph{all} these Shilov boundaries determine the
actual weight, including all poles.

Fix a normalized energy sublevel
$\G_C=\{u\in\G:\sup u=0,\ E(u)\ge-C\}$.
Its evaluation map on $\mathscr S$ is injective and takes values in
\[
 \prod_{\xi\in\mathscr S}[-L_\xi(C),0],\qquad
 L_\xi(C)=\frac{E^*(\MA(w_\xi))+C}{\alpha_\xi}.
\]
Tychonoff compactness implies that every net of such evaluations
has a coordinatewise convergent subnet with finite coordinates.
This is compactness in the product space of boundary values.
It does \emph{not} show that the limit coordinates are represented
by an actual member of $\G$, that the normalized image is closed,
or that energy is semicontinuous there. Those are still mathematical
gates for actual energy-sublevel compactness.
Likewise this point bound alone does not prove uniqueness from
equality of generated MA measures.

\begin{remark}
We have not claimed that $\delta_x$ itself has finite full bounded
dual merely because it is part of $\MA(w_x)$.
The whole test measure has the supporting-energy formula
\eqref{dual}; its positivity is exactly what gives the atomic
point inequality. The distinction retains the raw mass-two dual
and avoids an invalid componentwise assertion.
\end{remark}

\section{Source scope and original residual}
The exact CLDv2 full source is bound locally by SHA256
\texttt{64846cb00926267c6238d1afac61b5eaed0a0e3a6349725738f9b1bf163de085}.
Applied immediate readings include 1.7.1 and1.7.6--8
(LF1773--1830 and1931--2135), the complete local finite-flat block
11.6.1--5 (LF15846--15912), the graded Abhyankar criterion13.8.1--5
(LF19772--19820), the full calibration formula/proof14.5.2--3
(LF20953--21008), and the entire atomic calculation8.9.1--4
(LF10595--10739). The LF windows actually retained may include
adjacent context; essential statements and whole immediate proofs
were read in bounded outputs. The full relative atom formula18.6.2
(LF27607--27671) was read as corroboration, not a premise of the
finite-flat/calibration proof.
No algebraically closed standing hypothesis from14.6 is imported.
The real-plane coefficient one was computed above rather than
inferred from an OCR inequality.

The continuous products and bounded Green foundations retain the
exact earlier 8.5--8.8 source applications and all their hypotheses;
the complete fresh point construction and energy argument are
given here. GR is the full17September2025 v2 source, bound by
\texttt{3a221b86116604c950db07e4c60ab184701f0624deebabd92502e5c81871436b}.
Lemma6.3 and its whole immediate proof were read at LF1673--1780.
Its deeper Temkin/Ducros reduction proofs, basic Berkovich
relative-interior/dimension and completed-residue facts, graded
field transcendence, and the analytic quotient construction are
disclosed conventional foundations. There is no claim of an
entire new recursive source audit. All full primary bodies,
literal units, pixels and book contexts remain LOCALONLY.

Original P16.7 still requires actual separated energy-sublevel
compactness/representability, envelopes or a proved variational
replacement, arbitrary-support equation existence at the stated
full bounded-dual criterion, and full generated pointwise uniqueness.
This paper proves the previously missing actual point-finiteness
and coordinate lower bounds. It neither closes the original nor
changes any accepted count. The protected P16.4 editor remains open
and unchanged; new work is not sent as hot feedback to active
native campaigns. All claims remain UNREVIEWED.

\end{document}
